%=====================================================================
\chapter{Rule-Based Expert Systems}\label{ch:expert}
%=====================================================================
\locator{4}{Part IV --- Knowledge, Logic and Reasoning}

\begin{outcomes}
By the end of this chapter you will be able to:
\begin{tight}
  \item \textbf{Build} a production system: working memory, a rule base and an
        inference engine running the match--resolve--act cycle.
  \item \textbf{Apply} forward and backward chaining, and \textbf{choose} between
        them for a given task.
  \item \textbf{Resolve} conflicts between applicable rules, and \textbf{justify}
        the strategy chosen.
  \item \textbf{Produce} an explanation of a conclusion, and \textbf{argue} why it
        is the artefact an auditor accepts.
  \item \textbf{Describe} the knowledge acquisition bottleneck and brittleness, and
        \textbf{say} what they did to the expert-system industry.
\end{tight}
\end{outcomes}

\begin{prereq}
\chref{propositional} and \chref{fol} (inference, Horn clauses, chaining as
resolution under a restriction) and \chref{kr} (the representation the rules act
on). This chapter is where Part~IV becomes a product.
\end{prereq}

\begin{keyterms}
production rule $\bullet$ working memory $\bullet$ rule base $\bullet$ inference
engine $\bullet$ match--resolve--act $\bullet$ forward chaining $\bullet$ backward
chaining $\bullet$ conflict set $\bullet$ conflict resolution $\bullet$ specificity
$\bullet$ recency $\bullet$ refraction $\bullet$ explanation $\bullet$ knowledge
acquisition bottleneck $\bullet$ brittleness $\bullet$ graceful degradation
\end{keyterms}

\section{The Escalation Adviser}

\chref{agents} introduced it and deferred it: at 21:00 the queue holds three hundred
open tickets, one engineer is on call, and something must decide what to do with
each one. The decision depends on severity, customer tier, whether there is a
workaround, how much of the service level agreement is left, whether the service
owner is reachable, and how often the on-call engineer has already been woken.

Nobody can write that as a formula, and everybody in the operations team can state
the individual rules. That combination --- many small, separately intelligible
conditions and no closed form --- is exactly what a \term{rule-based expert system}
is for, and it is still the commonest deployed form of symbolic AI.

\begin{definitionbox}[Production system]
A \term{production system} has three parts:
\begin{tight}
  \item a \term{working memory} of facts currently believed;
  \item a \term{rule base} of \term{production rules}, each a condition--action
        pair, conventionally written \textsc{if} conditions \textsc{then}
        conclusion;
  \item an \term{inference engine} that repeatedly runs the
        \term{match--resolve--act} cycle: find every rule whose conditions hold
        (the \term{conflict set}), choose one, and apply its conclusion to working
        memory.
\end{tight}
The rules do not call each other. Each fires when the facts happen to satisfy it,
so a rule can be added or corrected without anyone tracing control flow --- which
is the property that lets a domain expert maintain the system.
\end{definitionbox}

\dsfig{%
\begin{tikzpicture}[font=\scriptsize,
  box/.style={rounded corners=3pt, draw=#1, line width=1.1pt, fill=#1!8,
              text=#1!50!black, text width=3.5cm, align=flush left,
              inner sep=6pt, minimum height=2.3cm},
  hd/.style={rounded corners=2pt, fill=#1, text=white,
             font=\scriptsize\bfseries, inner sep=3pt, text width=3.3cm,
             align=center},
  eng/.style={rounded corners=3pt, draw=primary, line width=1.4pt,
              fill=primary!12, text=primary!65!black,
              font=\scriptsize\bfseries, minimum width=2.5cm,
              minimum height=0.62cm, align=center},
  a/.style={-{Stealth[length=2.2mm]}, neutral!75, line width=1pt},
  cyc/.style={-{Stealth[length=2mm]}, primary!60, line width=1pt}]

\node[box=secondary] (wm) at (-4.9,0)
  {\\[2pt]\texttt{severity = critical}\\
   \texttt{tier = platinum}\\
   \texttt{service = payments}\\
   \texttt{workaround = no}\\
   \texttt{owner\_available = no}\\[3pt]
   {\itshape and everything derived, added as it is derived}};

\node[box=goldhl] (rb) at (4.9,0)
  {\\[2pt]\texttt{R1: severity=critical}\\
   \quad\texttt{-> priority=P1}\\[2pt]
   \texttt{R5: priority=P1}\\
   \quad\texttt{-> needs\_owner=yes}\\[3pt]
   {\itshape 20 working rules, each separately readable and
   separately correctable --- plus one deliberately bad one}};

\node[hd=secondary] at (wm.north) {WORKING MEMORY};
\node[hd=goldhl]    at (rb.north) {RULE BASE};

\node[eng] (m) at (0,0.85)  {1. MATCH};
\node[eng] (r) at (0,0.0)   {2. RESOLVE};
\node[eng] (c) at (0,-0.85) {3. ACT};
\node[font=\tiny\itshape, text=neutral, align=center, text width=3.0cm]
  at (0,-1.7) {the inference engine, until no rule is applicable};

\draw[a] (wm.east) -- (m.west);
\draw[a] (rb.west) -- (m.east);
\draw[cyc] (m) -- (r);
\draw[cyc] (r) -- (c);
\draw[a] (c.west) .. controls +(-1.6,0) and +(0,-1.6) .. (wm.south)
  node[midway, below, font=\tiny\itshape, text=neutral] {assert the conclusion};

\node[font=\tiny, text=neutral, anchor=west, align=left] at (0.95,0.85)
  {which rules apply now?};
\node[font=\tiny, text=neutral, anchor=west, align=left] at (0.95,0.0)
  {several do --- which first?};
\node[font=\tiny, text=neutral, anchor=west, align=left] at (0.95,-0.85)
  {add its conclusion};
\end{tikzpicture}}%
{The architecture. Working memory grows as conclusions are asserted, which makes
further rules applicable, and the cycle repeats until nothing new can be
derived.}{rule-engine}

\section{Forward Chaining}

Start from the facts and derive everything that follows. This is
\term{forward chaining}, and it is \emph{data-driven}: the arrival of a new fact
triggers whatever it triggers.

On ticket T-4471 --- critical, platinum, payments, under two hours of service level
agreement left, owner unreachable, on-call budget available --- the engine fires
nine rules in ten cycles, testing $188$ conditions, and reaches
\texttt{action = wake\_on\_call}. Along the way it also derives that the account
team must be notified, that finance must be notified because the service is
payments, that the payments incident runbook applies, and that the page is urgent
--- none of which anybody asked about.

That is the character of forward chaining. It is the right choice for monitoring
and alerting, where you want every consequence of an arriving fact. It is the wrong
choice when you have one question, because it derives a great deal you did not ask
for.

\section{Backward Chaining}

Start from the question and work backwards to the facts. To establish a goal, find a
rule concluding it and recursively establish that rule's conditions.

Asked only \emph{why is T-4471 escalated?}, backward chaining visits six goals
against forward chaining's nine firings, and --- far more importantly --- it
produces a proof tree rather than a set of conclusions.

\dsfig{%
\begin{tikzpicture}[font=\scriptsize,
  g/.style={rounded corners=3pt, draw=primary, line width=1pt, fill=primary!10,
            text=primary!60!black, font=\ttfamily\scriptsize,
            align=center, inner sep=4pt},
  f/.style={rounded corners=3pt, draw=greenhl, line width=1pt, fill=greenhl!12,
            text=greenhl!45!black, font=\ttfamily\scriptsize,
            align=center, inner sep=4pt},
  e/.style={-{Stealth[length=1.8mm]}, neutral!70, line width=0.9pt},
  rl/.style={font=\tiny\bfseries, text=accent, inner sep=1.5pt}]

\node[g] (esc)  at (0,2.6)    {escalate = yes};
\node[g] (need) at (-2.5,1.3) {needs\_owner = yes};
\node[f] (own)  at (2.4,1.3)  {owner\_available = no};
\node[g] (pri)  at (-2.5,0.0) {priority = P1};
\node[f] (sev)  at (-2.5,-1.3){severity = critical};

\draw[e] (esc) -- node[rl, above left] {R9} (need);
\draw[e] (esc) -- node[rl, above right] {R9} (own);
\draw[e] (need) -- node[rl, left] {R5} (pri);
\draw[e] (pri) -- node[rl, left] {R1} (sev);

\node[font=\tiny\itshape, text=greenhl!45!black, anchor=west] at (3.6,1.3)
  {given in the ticket};
\node[font=\tiny\itshape, text=greenhl!45!black, anchor=west] at (-1.5,-1.3)
  {given in the ticket};

\node[rounded corners=3pt, fill=lightbg, draw=primary!30, line width=0.9pt,
      text=primary!70!black, font=\scriptsize, text width=13.0cm,
      align=flush left, inner sep=6pt] at (0,-2.7)
  {\textbf{Read it upward and it is an argument in English:} the severity is
   critical, so by R1 the priority is P1; a P1 needs its service owner, by R5; the
   owner is not available; therefore by R9 the ticket is escalated.

   \smallskip
   That is the artefact. It names every rule and every fact it relied on, a person
   can check each line, and if the conclusion is wrong it identifies exactly which
   rule to correct. No probability and no learned model produces this, and
   \chref{responsible} argues it is an accountability instrument rather than a
   convenience.};
\end{tikzpicture}}%
{The proof tree for \emph{why is T-4471 escalated?} Rounded green boxes are facts
from the ticket; blue boxes are derived, each labelled with the rule that derived
it.}{explanation}

\smallskip
\noindent\begin{longtable}{@{}L{3.0cm}L{5.8cm}L{6.0cm}@{}}
\toprule
\rowcolor{primary!15}
 & \textbf{Forward chaining} & \textbf{Backward chaining} \\
\midrule
\endhead
Driven by & the data: a new fact triggers rules & the goal: a question triggers
subgoals \\
Derives & everything that follows & only what bears on the question \\
On T-4471 & 9 rules fired, 188 condition tests & 6 goals visited \\
Natural output & a set of conclusions & a proof tree \\
Good for & monitoring, alerting, updating a dashboard & diagnosis, advice,
answering ``why?'' \\
Wasteful when & you have one specific question & you want every consequence \\
\bottomrule
\caption{The two directions. Most deployed systems use both: forward chaining to
maintain state, backward chaining to answer questions about
it.}\label{tab:chaining}
\end{longtable}

\section{Conflict Resolution}

At most cycles several rules apply at once, and the order changes the outcome ---
sometimes only in efficiency, sometimes in the conclusion. Choosing is \term{conflict
resolution}, and the standard strategies are:

\begin{tight}
  \item \term{Specificity} --- prefer the rule with more conditions. A rule that
        mentions \emph{major severity and platinum tier} is about a narrower
        situation than one mentioning \emph{major severity} alone, and the narrower
        rule is the one the expert meant to apply. This is the workhorse.
  \item \term{Recency} --- prefer the rule matching the most recently added facts,
        which keeps the system on one line of reasoning instead of drifting.
  \item \term{Refraction} --- never fire the same rule on the same facts twice. Not
        an optimisation: without it the engine loops forever.
  \item \term{Explicit priority} --- a numeric salience on each rule. Always
        available, and a warning sign: if a rule base needs many priorities, the
        rules are probably overlapping in ways nobody has thought through.
\end{tight}

\begin{alertbox}[Specificity is worth a night's sleep]
The rule base contains one deliberately over-general rule, R0: \emph{a major
ticket is P1}. It is listed first.

Ticket T-4472 is major severity on a \emph{standard}-tier customer with a
workaround. Under plain rule-order resolution the engine reaches R0 first, makes
the ticket P1, and ends at \texttt{wake\_on\_call} --- an engineer woken at three in
the morning for a ticket that has a workaround. Under specificity the
two-condition R3 is preferred, the ticket is P2, and the action is
\texttt{leave\_in\_queue}.

Same rules, same facts, different conflict-resolution strategy, and the difference
is somebody's night. That is why a rule base is not a list of \textsc{if}
statements: a list has the order whoever typed it chose, while specificity derives
the order from how specific each rule's claim is --- which is a property of the
knowledge rather than of the file.
\end{alertbox}

\section{What Conflict Resolution Does Not Fix}

Now run the two directions against each other on that same ticket, and something
uncomfortable happens.

Forward chaining, with specificity, concludes \texttt{escalate = no}. Backward
chaining, asked to prove \texttt{escalate = yes}, \emph{succeeds} --- via R5 from
$\mathit{priority} = \text{P1}$, which it obtains from R0. The two directions
disagree about the same ticket in the same rule base.

\begin{conceptbox}
Backward chaining has no conflict set. It asks ``is there \emph{any} rule
concluding this goal?'', and R0 is one. Forward chaining never used R0 only because
specificity reached R3 first.

So specificity was not repairing the rule base; it was \emph{masking} a bad rule.
R0 asserts that every major ticket is P1, which is simply false, and no
conflict-resolution strategy makes a false rule true --- it merely arranges for
something else to be consulted first, on the inputs you happened to test.

The fix is to delete R0, whose work is already done correctly by R2 and R3. With it
gone the two directions agree, and the lab checks that they do.

\smallskip
The general lesson is worth carrying beyond rule engines: \emph{when two correct
procedures give different answers from the same knowledge, the knowledge is
inconsistent}. Do not tune the procedure. And note how the inconsistency was found
--- not by inspection, but by asking the same question two ways and comparing.
That is a cheap and underused test for any rule base.
\end{conceptbox}

\section{Why the Industry Collapsed}

Rule engines work. MYCIN diagnosed bacterial infections in the 1970s at the level of
faculty specialists; XCON configured VAX systems and saved Digital tens of millions
of dollars a year. By the mid-1980s there was a substantial commercial industry, and
by 1993 it had collapsed. Two causes, and both are still live.

\textbf{The knowledge acquisition bottleneck.} Every rule has to be elicited from an
expert, by an engineer, one at a time. Experts are busy, and --- worse --- they are
frequently unable to state what they do: the knowledge is compiled into judgement.
A thousand-rule system represents years of interviews and cannot be assembled any
faster. Meanwhile, nothing in the system learns; it fails
\chref{introduction}'s fourth test and it will still be failing it next year.

\textbf{Brittleness.} A rule base is confident inside its competence and gives no
signal at the edge. Present a ticket whose situation nobody anticipated and the
system does not degrade gracefully --- it either fires nothing, or fires a rule
whose conditions technically match and whose conclusion is absurd. A human
specialist says ``I have not seen this before''; a rule base has no representation
of its own boundary.

\begin{conceptbox}
The historical lesson is the one from \chref{introduction}: what failed was a claim
about \emph{reach}, not the method. Rule engines are running in production today in
every bank, insurer and operations centre, generally under a name like ``business
rules engine'' or ``decision tables'', and generally with exactly the architecture of
\dsref{rule-engine}.

And the two failures are precisely the two strengths of what came next. Machine
learning acquires its knowledge from data instead of from interviews, which dissolves
the bottleneck; and it degrades smoothly rather than falling off a cliff. It gives up
the explanation tree of \dsref{explanation} to do so. That trade --- explicability
for scale --- is the central tension of \chref{responsible}, and it is worth
noticing that this course has now built both sides of the argument.
\end{conceptbox}

\begin{worked}[One ticket, two budgets, two different actions]
Tickets T-4471 and T-4473 differ in only one fact, and the difference is not about
the ticket at all.

\smallskip
\textbf{T-4471:} critical, platinum, payments, under two hours left, owner
unavailable, \texttt{wake\_budget = available}.

\smallskip
The engine fires, in order: R17 (critical and platinum, so notify the account team),
R1 (critical, so P1), R18 (payments and P1, so notify finance), R5 (P1, so needs the
owner), R9 (needs the owner, who is unavailable, so escalate), R15 (escalate with
budget available, so \emph{wake the on-call engineer}), R19 (escalate on payments,
so the payments runbook applies), R20 (waking someone for a platinum customer, so
the page is urgent), R12 (under two hours, so the risk is high). Nine rules, and
the answer is \texttt{wake\_on\_call}.

\smallskip
\textbf{T-4473:} major, platinum, auth, over two hours left, owner unavailable,
\texttt{wake\_budget = spent}.

\smallskip
R2 (major and platinum, so P1 --- note the two-condition rule, chosen by
specificity over the one-condition major rule), R5 (P1, so needs the owner), R9
(escalate), R14 (escalate but \emph{the budget is spent}, so
\texttt{queue\_for\_morning}), R13 (over two hours, so risk is low). Five rules, and
the answer is \texttt{queue\_for\_morning}.

\smallskip
\textbf{What to notice.} Both tickets are escalated --- \texttt{escalate = yes} in
both working memories. They receive different \emph{actions} because R14 and R15
consult a fact about the \emph{engineer}, not about the ticket. That is exactly the
internal state \chref{agents} showed a model-based agent must maintain, and here it
is one more fact in working memory and one more pair of rules.

\smallskip
\textbf{And notice what it costs to change.} Suppose the operations lead decides
that platinum customers may wake the engineer even past the budget. That is one new
rule, with three conditions so specificity puts it ahead of R14, and \emph{no
existing rule is touched}. Compare making the same change inside nested
\texttt{if} statements, where the modification is entangled with everything around
it and nobody can be sure what else moved.
\end{worked}

\begin{pitfall}
\begin{tight}
  \item \textbf{Omitting refraction.} The engine re-fires the same rule on the same
        facts forever. The symptom is a hang on the first realistic input.
  \item \textbf{Relying on rule order.} If the answer changes when the rules are
        rearranged, the rule base is under-specified; add the missing condition
        rather than reordering.
  \item \textbf{Forward chaining to answer one question.} It derived nine
        conclusions where six goals would have done, and none of the extras was
        asked for.
  \item \textbf{Treating explanation as a feature to add later.} The proof tree is
        a by-product of backward chaining and nearly impossible to bolt on
        afterwards. It is also the reason the system will be allowed into
        production.
  \item \textbf{Many explicit priorities.} A salience number is a patch over
        overlapping rules. More than a handful means the rule base needs
        restructuring, not more numbers.
  \item \textbf{Expecting graceful degradation.} At the edge of its competence a
        rule base is confidently wrong rather than uncertain, and it will not tell
        you which.
\end{tight}
\end{pitfall}

%---------------------------------------------------------------------
\begin{lab}[A Twenty-Rule Escalation Shell That Explains Itself]
Forward chaining with conflict resolution, backward chaining with an explanation
tree, and the measurement of what each costs.
\begin{lstlisting}[language=Python]
"""Chapter 13 lab -- a rule-based escalation adviser.

Twenty rules an operations lead could read and correct. The engine
fires until nothing new follows; the explanation tree is what makes
the result auditable, and it is a by-product of backward chaining
rather than something added afterwards.
"""

# (name, [conditions], conclusion); each is an (attribute, value) pair
RULES = [
    # R0 is a deliberately general fallback, and it is listed FIRST.
    # Plain rule-order resolution therefore reaches it before the
    # narrower R2 and R3, and gets the wrong answer. Specificity does
    # not, which is the point of the comparison below.
    ("R0",  [("severity", "major")], ("priority", "P1")),
    ("R1",  [("severity", "critical")], ("priority", "P1")),
    ("R2",  [("severity", "major"), ("tier", "platinum")],
     ("priority", "P1")),
    ("R3",  [("severity", "major"), ("tier", "standard")],
     ("priority", "P2")),
    ("R4",  [("severity", "minor")], ("priority", "P3")),
    ("R5",  [("priority", "P1")], ("needs_owner", "yes")),
    ("R6",  [("priority", "P2"), ("workaround", "no")],
     ("needs_owner", "yes")),
    ("R7",  [("priority", "P2"), ("workaround", "yes")],
     ("needs_owner", "no")),
    ("R8",  [("priority", "P3")], ("needs_owner", "no")),
    ("R9",  [("needs_owner", "yes"), ("owner_available", "no")],
     ("escalate", "yes")),
    ("R10", [("priority", "P1"), ("sla_risk", "high")],
     ("escalate", "yes")),
    ("R11", [("needs_owner", "no")], ("escalate", "no")),
    ("R12", [("sla_hours_left", "under2")], ("sla_risk", "high")),
    ("R13", [("sla_hours_left", "over2")], ("sla_risk", "low")),
    ("R14", [("escalate", "yes"), ("wake_budget", "spent")],
     ("action", "queue_for_morning")),
    ("R15", [("escalate", "yes"), ("wake_budget", "available")],
     ("action", "wake_on_call")),
    ("R16", [("escalate", "no")], ("action", "leave_in_queue")),
    ("R17", [("severity", "critical"), ("tier", "platinum")],
     ("notify_account_team", "yes")),
    ("R18", [("service", "payments"), ("priority", "P1")],
     ("notify_finance", "yes")),
    ("R19", [("escalate", "yes"), ("service", "payments")],
     ("runbook", "payments_incident")),
    ("R20", [("action", "wake_on_call"), ("tier", "platinum")],
     ("page_severity", "urgent")),
]

TICKETS = {
    "T-4471": {"severity": "critical", "tier": "platinum",
               "service": "payments", "workaround": "no",
               "sla_hours_left": "under2", "owner_available": "no",
               "wake_budget": "available"},
    "T-4472": {"severity": "major", "tier": "standard",
               "service": "search", "workaround": "yes",
               "sla_hours_left": "over2", "owner_available": "no",
               "wake_budget": "available"},
    "T-4473": {"severity": "major", "tier": "platinum",
               "service": "auth", "workaround": "no",
               "sla_hours_left": "over2", "owner_available": "no",
               "wake_budget": "spent"},
}


class Trace:
    def __init__(self):
        self.fired = []
        self.tests = 0
        self.cycles = 0


def matches(rule, wm, t):
    for (attr, val) in rule[1]:
        t.tests += 1
        if wm.get(attr) != val:
            return False
    return True


def forward_chain(facts, resolve="specificity", rules=None):
    """Match, resolve, act -- until nothing new follows.

    The 'fired' set is REFRACTION: without it the engine re-fires the
    same rule on the same facts forever.

    'r[2][0] not in wm' means each attribute is DECIDED ONCE, by the
    rule conflict resolution picks. Without it a less specific rule
    that fires later would quietly overwrite a more specific rule's
    conclusion, and conflict resolution would be pointless.
    """
    rules = RULES if rules is None else rules
    wm, t, fired = dict(facts), Trace(), set()
    while True:
        t.cycles += 1
        conflict = [r for r in rules
                    if r[0] not in fired
                    and matches(r, wm, t)
                    and r[2][0] not in wm]
        if not conflict:
            return wm, t
        if resolve == "specificity":
            # more conditions first; ties broken by rule order
            conflict.sort(key=lambda r: (-len(r[1]), rules.index(r)))
        else:
            conflict.sort(key=lambda r: rules.index(r))
        r = conflict[0]
        wm[r[2][0]] = r[2][1]
        fired.add(r[0])
        t.fired.append(r[0])


def backward_chain(goal, facts, seen=None):
    """Prove an (attribute, value) goal. Returns a proof tree."""
    if seen is None:
        seen = set()
    if facts.get(goal[0]) == goal[1]:
        return ("fact", goal)
    if goal in seen:
        return None                 # a cycle in the rule base
    seen = seen | {goal}
    for r in RULES:
        if r[2] != goal:
            continue
        subs = []
        for c in r[1]:
            got = backward_chain(c, facts, seen)
            if got is None:
                subs = None
                break
            subs.append(got)
        if subs is not None:
            return ("rule", r[0], goal, subs)
    return None


def explain(tree, indent=0):
    pad = "  " * indent
    if tree[0] == "fact":
        return ["%s%s = %s   (given in the ticket)"
                % (pad, tree[1][0], tree[1][1])]
    _, name, goal, subs = tree
    out = ["%s%s = %s   (by %s, because)" % (pad, goal[0], goal[1], name)]
    for s in subs:
        out += explain(s, indent + 1)
    return out


def goals_visited(goal, facts, seen=None, count=None):
    """How much work backward chaining does for ONE question."""
    if count is None:
        count = [0]
    count[0] += 1
    if seen is None:
        seen = set()
    if facts.get(goal[0]) == goal[1]:
        return True, count[0]
    if goal in seen:
        return False, count[0]
    seen = seen | {goal}
    for r in RULES:
        if r[2] != goal:
            continue
        if all(goals_visited(c, facts, seen, count)[0] for c in r[1]):
            return True, count[0]
    return False, count[0]


if __name__ == "__main__":
    print("rule base: %d rules" % len(RULES))
    print()
    print("ticket    action              rules  tests  cycles")
    print("-" * 52)
    results = {}
    for tid, facts in TICKETS.items():
        wm, t = forward_chain(facts)
        results[tid] = (wm, t)
        print("%-9s %-19s %5d %6d %7d"
              % (tid, wm.get("action"), len(t.fired), t.tests, t.cycles))

    # the three actions differ, and only one fact separates the last two
    assert results["T-4471"][0]["action"] == "wake_on_call"
    assert results["T-4472"][0]["action"] == "leave_in_queue"
    assert results["T-4473"][0]["action"] == "queue_for_morning"
    print()
    print("T-4471 firing order:", " ".join(results["T-4471"][1].fired))
    print("T-4473 firing order:", " ".join(results["T-4473"][1].fired))
    print()
    print("Both are escalated; they differ only in wake_budget, and")
    print("that is a fact about the ENGINEER, not about the ticket.")
    assert results["T-4471"][0]["escalate"] == "yes"
    assert results["T-4473"][0]["escalate"] == "yes"

    # --- specificity is doing real work --------------------------
    print()
    print("=== specificity against plain rule order ===")
    print("ticket    specificity         rule order")
    print("-" * 46)
    for tid, facts in TICKETS.items():
        a, _ = forward_chain(facts, "specificity")
        b, _ = forward_chain(facts, "order")
        flag = "  <-- DIFFER" if a["action"] != b["action"] else ""
        print("%-9s %-19s %-19s%s"
              % (tid, a["action"], b["action"], flag))

    spec, _ = forward_chain(TICKETS["T-4472"], "specificity")
    order, _ = forward_chain(TICKETS["T-4472"], "order")
    assert spec["priority"] == "P2" and spec["action"] == "leave_in_queue"
    assert order["priority"] == "P1" and order["action"] == "wake_on_call"
    print()
    print("T-4472 is major severity on a STANDARD tier customer.")
    print("R0 ('major -> P1') is general and listed first, so plain")
    print("rule order reaches it and wakes an engineer at 3 a.m. for")
    print("a standard-tier ticket with a workaround. Specificity")
    print("prefers the two-condition R3, gets P2, and leaves it in")
    print("the queue. That is a real engineer's night, decided by the")
    print("conflict-resolution strategy alone.")

    # --- the explanation ------------------------------------------
    print()
    print("=== why is T-4471 escalated? ===")
    tree = backward_chain(("escalate", "yes"), TICKETS["T-4471"])
    assert tree is not None
    for line in explain(tree):
        print("   " + line)

    # --- the two directions DISAGREE, and that is a real bug -----
    print()
    print("=== and is T-4472 escalated? ===")
    fwd_wm, _ = forward_chain(TICKETS["T-4472"])
    up = backward_chain(("escalate", "yes"), TICKETS["T-4472"])
    print("   forward chaining says escalate =", fwd_wm["escalate"])
    print("   backward chaining proves escalate = yes:", up is not None)
    assert fwd_wm["escalate"] == "no" and up is not None
    print()
    print("   They disagree. Backward chaining's proof:")
    for line in explain(up):
        print("      " + line)
    print()
    print("   Backward chaining has no conflict resolution: it asks")
    print("   'is there ANY rule concluding this?' and R0 is one.")
    print("   Forward chaining never used R0 because specificity")
    print("   preferred R3 -- so specificity was MASKING a bad rule,")
    print("   not fixing it. R0 says every major ticket is P1, which")
    print("   is simply false, and no resolution strategy makes a")
    print("   false rule true.")

    # the fix is to delete the over-general rule, not to tune the
    # conflict-resolution strategy
    print()
    print("=== delete R0 and ask both directions again ===")
    fixed = [r for r in RULES if r[0] != "R0"]
    fwd2, _ = forward_chain(TICKETS["T-4472"], rules=fixed)
    globals()["RULES"], saved = fixed, RULES
    up2 = backward_chain(("escalate", "yes"), TICKETS["T-4472"])
    down2 = backward_chain(("escalate", "no"), TICKETS["T-4472"])
    globals()["RULES"] = saved
    print("   forward says escalate =", fwd2["escalate"])
    print("   backward proves escalate = yes:", up2 is not None)
    assert fwd2["escalate"] == "no" and up2 is None and down2 is not None
    print("   They now agree. The explanation:")
    for line in explain(down2):
        print("      " + line)

    # --- what each direction costs --------------------------------
    fwd = len(results["T-4471"][1].fired)
    _, bwd = goals_visited(("escalate", "yes"), TICKETS["T-4471"])
    print()
    print("forward chaining : %d rules fired to derive everything" % fwd)
    print("backward chaining: %d goals visited to answer one question"
          % bwd)
    assert bwd < fwd
    print()
    print("Forward chaining also concluded notify_finance, runbook and")
    print("page_severity -- useful for a dashboard, wasted for a")
    print("single question.")

    # --- adding a rule changes nothing else -----------------------
    print()
    print("=== the operations lead changes the policy ===")
    extended = RULES + [("R21", [("escalate", "yes"),
                                 ("wake_budget", "spent"),
                                 ("tier", "platinum")],
                         ("action", "wake_on_call"))]
    wm_new, t_new = forward_chain(TICKETS["T-4473"], rules=extended)
    print("   'platinum may wake the engineer past the budget'")
    print("   T-4473 action was queue_for_morning, now:",
          wm_new["action"])
    print("   firing order:", " ".join(t_new.fired))
    assert wm_new["action"] == "wake_on_call"
    print("   One rule added. No existing rule was touched, and")
    print("   specificity put it ahead of R14 automatically.")
\end{lstlisting}
Then remove the \texttt{fired} set from \texttt{forward\_chain} --- that is
refraction --- and re-run. The engine never terminates, because R1 is applicable
forever once the severity is critical. It is worth seeing once.
\end{lab}

%---------------------------------------------------------------------
\begin{checkpoint}
\begin{tightnum}
  \item T-4471 and T-4473 are both escalated and receive different actions. Name the
        single fact that separates them, say whether it is a fact about the ticket,
        and relate it to the agent architectures of \chref{agents}.
  \item Forward chaining fired nine rules on T-4471; backward chaining answered
        ``why escalated?'' in six goals. Give one task for which the nine is the
        right answer and the five is the wrong one.
  \item An auditor asks why a particular customer's ticket was left in the queue
        overnight. Explain what your system hands them, and why a model that merely
        output a probability would not satisfy the request.
\end{tightnum}
\end{checkpoint}

%---------------------------------------------------------------------
\begin{chaptersummary}
\begin{tight}
  \item A \term{production system} is working memory, a rule base and an inference
        engine running \term{match--resolve--act}. Rules do not call each other, so
        one can be added or corrected without tracing control flow --- which is what
        lets a domain expert maintain it.
  \item \term{Forward chaining} is data-driven and derives everything that follows:
        nine rules and $188$ condition tests on one ticket. Right for monitoring;
        wasteful for a single question.
  \item \term{Backward chaining} is goal-driven, visits only what bears on the
        question --- six goals here --- and produces a \term{proof tree}.
  \item \term{Conflict resolution} chooses among applicable rules.
        \term{Specificity} is the workhorse and encodes a real property of the
        knowledge; \term{refraction} is not an optimisation but the thing that makes
        the engine terminate.
  \item The \term{explanation} is the deliverable. It names every rule and fact the
        conclusion rests on, a person can check it line by line, and it identifies
        which rule to correct when the answer is wrong.
  \item Conflict resolution can \emph{mask} a bad rule without fixing it. Asking the
        same question forwards and backwards made the two directions disagree on
        T-4472 and exposed an over-general rule that specificity had been hiding.
        When two correct procedures differ on the same knowledge, the knowledge is
        inconsistent --- fix the rules, not the strategy.
  \item The industry collapsed on the \term{knowledge acquisition bottleneck} ---
        rules elicited one at a time from busy experts, and nothing learns --- and
        on \term{brittleness}, confident wrongness at the edge of competence.
  \item Both failures are the strengths of machine learning, which acquires
        knowledge from data and degrades smoothly, and which gives up the
        explanation tree to do it. That trade is the subject of
        \chref{responsible}.
\end{tight}
\end{chaptersummary}

\begin{reviewq}
  \item \textbf{(MCQ)} The conflict set is: (a)~rules that contradict each other
        (b)~rules whose conditions currently hold (c)~facts that disagree
        (d)~goals not yet proved.
  \item \textbf{(MCQ)} Refraction prevents: (a)~two rules firing at once
        (b)~a rule firing twice on the same facts (c)~backward chaining looping
        (d)~conflicting conclusions.
  \item \textbf{(MCQ)} Backward chaining is preferable when:
        (a)~a new fact has arrived and you want its consequences
        (b)~you have one specific question (c)~working memory is large
        (d)~the rules have many conditions.
  \item \textbf{(MCQ)} The knowledge acquisition bottleneck is:
        (a)~slow inference (b)~too many rules to store (c)~rules must be elicited
        from experts one at a time (d)~working memory overflow.
  \item \textbf{(Short)} Explain specificity as a conflict-resolution strategy and
        why it encodes something real about the knowledge rather than being an
        arbitrary tie-break.
  \item \textbf{(Short)} Why is an explanation tree difficult to add to a system
        afterwards, and easy to obtain from backward chaining?
  \item \textbf{(Short)} State the two causes of the expert-system collapse and,
        for each, say how machine learning addresses it and what is given up.
  \item \textbf{(Applied)} Write a rule implementing ``a P1 on a platinum customer
        with no workaround must notify the account team even if the owner is
        available''. Give its conditions and conclusion, say where specificity
        places it relative to R17, and trace what T-4473 would do with it added.
\end{reviewq}
